%Discussion section 
This section interprets the experimental results in terms of the main constraints on haptic performance: communication latency, delayed contact behaviour, hardware limitations, and the resulting design trade-offs.

\subsection{Communication and Latency Constraints} \label{sec:discussion_communication_latency}

The central finding from the communication evaluation is that transport delay, measured at 15.60--15.67\,ms mean round-trip with 95th and 99th percentiles around 16.2--16.4\,ms, is the dominant limitation on closed-loop performance. In contrast, host-side processing latency remained low (mean 0.299\,ms, 99th percentile 2.04\,ms), confirming that computation is not the bottleneck for the nominal 1\,kHz local haptic/device loops.

This delay directly bounds achievable haptic stiffness. As delay increases, phase lag reduces the stability margin and can inject artificial energy into the sampled-data loop \cite{colgatePassivityClassSampleddata1997}, so lower stiffness and additional damping are required. Consequently, the realised system renders compliant rather than rigid contact.

The ST-Link virtual COM interface was convenient for development, but the observed burst-like delivery pattern and low-rate packet loss suggest USB CDC buffering and host scheduling overhead that is large relative to the control timestep. Although the present work does not isolate every stage in the communication path, the evidence shows that communication architecture must be treated as a first-order design constraint alongside control and mechanical design.

The system should therefore be interpreted as a high-rate local haptic renderer operating over delayed sampled state, rather than as a true 1\,kHz end-to-end haptic loop. Snapshot channels reduce backlog-induced latency by ensuring that the haptic and device loops consume the newest available state, but they cannot remove transport delay, increase the physics-simulation update rate, or guarantee that the rendered contact state is physically fresh at 1\,ms intervals.

\subsection{Local Contact State Limitation} \label{sec:discussion_local_contact_state}

A further limitation of the realised architecture is that the nominal 1\,kHz haptic loop did not evolve a complete local contact state independently of the slower simulation and communication layers. It consumed latest-available snapshots of tool and world state, then recomputed proxy/contact information at each cycle. The renderer could therefore execute at high frequency, but the physical meaning of each update was still bounded by snapshot freshness and communication delay. The limitation was not the execution rate of the haptic computation, but that the rendered state was not itself evolved and validated at 1\,kHz across the full interaction loop.

A more complete architecture would maintain persistent proxy and local contact state inside the high-rate loop. Rather than re-projecting the proxy from scratch, the proxy would evolve continuously subject to contact constraints, preserving contact history, tangential motion, and release behaviour. This is closer to constraint-based proxy or god-object approaches \cite{zillesConstraintbasedGodobjectMethod1995,ruspiniHapticDisplayComplex1997,salisburyHapticRenderingSurfaces1997}. The simpler projection-based approach used here was sufficient to demonstrate bounded contact, but it limited contact smoothness and weakens any claim to true 1\,kHz end-to-end haptic interaction.

This clarifies the relationship between the haptic renderer and the physics engine. The physics engine provided the authoritative global scene state at a lower rate, while the haptic loop consumed snapshots of that state. A stronger version of this multi-rate design would allow the haptic loop to maintain a reduced local model of contacted objects at high frequency, while slower application, rendering, or physics subsystems provide broader scene updates and maintain global consistency, following the motivation of multilevel haptic rendering for complex environments \cite{ruspiniHapticDisplayComplex1997}.

\subsection{Observed Contact Behaviour Under Delay} \label{sec:discussion_stability}

Although communication timing was imperfect, the tested interactions remained bounded and controllable because snapshot channels used the newest available state rather than queued historical samples, thread separation avoided cross-subsystem blocking, and damping countered delayed feedback effects. Embedded torque rate limiting further prevented abrupt torque transitions that can excite oscillations. Watchdog activation data also show that communication interruptions were handled safely: 259 activation events were recorded, with a mean duration of 2\,ms and torque removed during activation (Section~\ref{sec:results}).

These mechanisms address the passivity concerns identified by Colgate and Schenkel \cite{colgatePassivityClassSampleddata1997} in a practical rather than formal way. The results provide evidence of bounded behaviour in the tested cases, not a passivity proof. Contact therefore remained controllable without divergent motion or sustained oscillation, but the compliant feel was still a direct consequence of delayed feedback; rigid contact cannot be achieved under the measured latency conditions by force-model tuning alone.

\subsection{Hardware Limitations} \label{sec:discussion_hardware_limitations}
The hardware results must be interpreted within clear prototype constraints. The most important was the absence of current sensing. The initial design targeted STM32 B-G431-ESC1 motor drivers, but integration with the selected magnetic encoders proved unreliable within the project timescale. SimpleFOC Mini modules provided a more robust integration path, but left torque control effectively open-loop with respect to current. As a result, commanded torque cannot be interpreted as calibrated physical torque, and the hardware evaluation should be read primarily as evidence of logged command consistency and qualitative interaction capability. This limitation does not affect the communication, timing, and bounded-contact results, which do not depend on absolute force calibration.

A second limitation was the GT2 belt reduction used to increase base-joint torque. The belt improved available torque but slipped under high load, even after adding an adjustable idler pulley, probably because 3D-printed pulleys did not provide manufactured tooth-profile precision. Actuator torque limits would still bound force output even without slip, directly constraining renderable stiffness.

A further limitation arose from the failure of one actuator during development. Final hardware validation therefore used single-joint actuation with two-joint sensing, so conclusions about full 2-DOF force rendering remain provisional. Overall, actuator capability and transmission design impose significant constraints alongside communication latency in low-cost haptic systems.

\subsection{Design Trade-Offs} \label{sec:discussion_design_tradeoffs}
The system therefore reflects a trade-off between bounded contact behaviour, responsiveness, and force fidelity. Damping and torque rate limiting avoided abrupt or divergent responses, but reduced perceived stiffness. The measured limiter of approximately 55~N\,m\,s$^{-1}$ implies an approximate 0.13~s rise time to reach the 7~N\,m torque range, so hard contacts feel gradual rather than instantaneous.

An additional perceptual limitation was observed at contact release: after sustained penetration into a virtual wall, the transition back to free motion could feel slightly jerky. This may be associated with abrupt proxy-state changes, lag in the filtered proxy velocity state, or delayed torque commands. The implemented approach therefore prioritised controllability over maximum stiffness, showing that high-fidelity force rendering cannot be improved by optimising one subsystem in isolation.

\subsection{Project Management and Development Process} \label{sec:discussion_project_management}

The project plan provided a useful structure, but the final Gantt chart in Appendix~\ref{app:gantt} reflects an adaptive process shaped by integration risk. The main management decision was to prioritise a working system-level prototype over ideal subsystem performance: driver integration difficulties led to SimpleFOC Mini modules, and later actuator failure shifted final validation from full 2-DOF force feedback to single-axis actuation with two-axis sensing. Risk management therefore focused on preserving the communication, simulation, safety, and bounded-contact evaluation objectives through staged testing, torque limiting, slew-rate limiting, and watchdog safety.
